\documentclass[tikz,border=8pt]{standalone}
\usepackage{ctex}
\usepackage{tikz}
\usetikzlibrary{arrows.meta,positioning,fit,calc}
\setCJKmainfont{DFBiaoKaiShu-B5}
\setCJKsansfont{DFBiaoKaiShu-B5}

\begin{document}
\begin{tikzpicture}[
    node distance=8mm,
    chapterbox/.style={
        draw,
        rounded corners=3pt,
        align=left,
        font=\small,
        inner sep=7pt,
        minimum height=11mm,
        text width=14.0cm
    },
    component/.style={
        draw,
        rounded corners=2pt,
        align=left,
        font=\small,
        inner sep=5pt,
        minimum height=11mm,
        text width=6.4cm
    },
    componentwide/.style={
        draw,
        rounded corners=2pt,
        align=left,
        font=\small,
        inner sep=5pt,
        minimum height=11mm,
        text width=14.0cm
    },
    groupbox/.style={
        draw,
        dashed,
        rounded corners=4pt,
        inner sep=10pt
    },
    grouptitle/.style={
        font=\footnotesize\bfseries,
        fill=white,
        inner sep=1pt
    },
    flowarrow/.style={-{Stealth[length=2.8mm]}, thick}
]
    % ── 第一章 緒論 ──────────────────────────────────────────────────────────
    \node[chapterbox] (ch1) {\textbf{第一章 緒論}\\
    研究背景與動機\\
    研究問題與目標\\
    研究方法與貢獻\\
    論文架構};

    % ── 第二章 文獻回顧 ──────────────────────────────────────────────────────
    \node[component, below=10mm of ch1, xshift=-3.7cm] (b1) {線性迴歸在報酬預測中的\\典型定位與限制};
    \node[component, right=10mm of b1] (b2) {樹狀結構模型在非線性預測與\\特徵工程中的典型定位與限制};
    \node[component, below=7mm of b1] (b3) {技術指標如何轉化為深度學習\\可用之特徵表示};
    \node[component, below=7mm of b2] (b4) {LSTM 等時間序列模型在\\金融預測中的典型做法與限制};
    \node[component, below=7mm of b3] (b5) {圖神經網路如何建模\\股票之間的結構性關係};
    \node[component, below=7mm of b4] (b6) {GAT 及其變形在股票預測／排序\\任務上的應用與限制};
    \node[componentwide, below=7mm of b5, xshift=3.7cm] (b7) {因子投資情境下常用的訊號評估指標\\IC ／ ICIR};
    \node[groupbox, fit=(b1)(b7)] (groupb) {};
    \node[grouptitle, anchor=west] at ([xshift=6pt,yshift=3pt]groupb.north west) {第二章 文獻回顧};

    % ── 第三章 實驗設計 ──────────────────────────────────────────────────────
    \node[component, below=10mm of groupb, xshift=-3.7cm] (c1) {§\,3.1\enspace 資料來源與預處理};
    \node[component, right=10mm of c1] (c2) {§\,3.2\enspace 特徵工程與圖結構};
    \node[component, below=7mm of c2] (c3) {§\,3.3\enspace 模型設計};
    \node[component, below=7mm of c1] (c4) {§\,3.4\enspace 損失函數設計};
    \node[componentwide, below=7mm of c4, xshift=3.7cm] (c5) {§\,3.5\enspace 評估指標與基準比較};
    \node[componentwide, below=7mm of c5] (c6) {§\,3.6\enspace 統計假說檢定（Newey-West HAC $t$ 檢定、Block Bootstrap ICIR CI、Diebold-Mariano 配對比較）};
    \node[groupbox, fit=(c1)(c6)] (groupc) {};
    \node[grouptitle, anchor=west] at ([xshift=6pt,yshift=3pt]groupc.north west) {第三章 實驗設計};

    % ── 第四章 實驗結果 ──────────────────────────────────────────────────────
    \node[componentwide, below=10mm of groupc] (d0) {§\,4.1\enspace 實驗概述與市場背景};
    \node[component, below=7mm of d0, xshift=-3.7cm] (d1) {§\,4.2\enspace 非圖基準模型之實證分析\\（Linear、XGBoost、LSTM、DNN）};
    \node[component, right=10mm of d1] (d2) {§\,4.3\enspace 圖拓樸結構之排序效益比較\\（GAT-industry、universe、TwoGraph）};
    \node[component, below=7mm of d1] (d3) {§\,4.4\enspace 中性化機制分析\\（GAT-IndNeutral、GAT-full）};
    \node[component, below=7mm of d2] (d4) {§\,4.5\enspace 訓練視窗長度之跨模型效應\\（短期 / 中期 / 長期視窗）};
    \node[componentwide, below=7mm of d3, xshift=3.7cm] (d5) {§\,4.6\enspace 綜合詮釋與研究限制};
    \node[groupbox, fit=(d0)(d5)] (groupd) {};
    \node[grouptitle, anchor=west] at ([xshift=6pt,yshift=3pt]groupd.north west) {第四章 實驗結果};

    % ── 第五章 結論與未來工作 ────────────────────────────────────────────────
    \node[chapterbox, below=10mm of groupd] (ch5) {\textbf{第五章 結論與未來工作}\\
    主要發現（訓練視窗、圖拓樸、中性化機制、IC 與 Sharpe 詮釋）\\
    未來研究方向（交易成本、動態圖結構、市場狀態分組）};

    % ── 主流程箭頭 ───────────────────────────────────────────────────────────
    \draw[flowarrow] (ch1.south)    -- (groupb.north);
    \draw[flowarrow] (groupb.south) -- (groupc.north);
    \draw[flowarrow] (groupc.south) -- (groupd.north);
    \draw[flowarrow] (groupd.south) -- (ch5.north);

    % ── 第三章 內部箭頭 ──────────────────────────────────────────────────────
    \draw[flowarrow] (c1.east) -- (c2.west);
    \draw[flowarrow] (c2.south) -- (c3.north);
    \draw[flowarrow] (c3.west) -- (c4.east);
    \draw[flowarrow] (c4.south) -- (c5.north -| c4.south);
    \draw[flowarrow] (c5.south) -- (c6.north);

    % ── 第四章 內部箭頭 ──────────────────────────────────────────────────────
    \draw[flowarrow] (d0.south) -- (d0.south |- d1.north) -| (d1.north);
    \draw[flowarrow] (d0.south) -- (d0.south |- d2.north) -| (d2.north);
    \draw[flowarrow] (d1.south) -- (d3.north);
    \draw[flowarrow] (d2.south) -- (d4.north);
    \draw[flowarrow] (d3.south) -- (d3.south |- d5.north) -| (d5.north -| d3.south);

\end{tikzpicture}
\end{document}
